%==============================================================================
%  WEEKLY AI BRIEF -- issue #09, week ending 20 October 2026
%  Engine: tectonic
%==============================================================================
\documentclass[11pt,a4paper]{article}

%------------------------------ packages -------------------------------------
\usepackage[T1]{fontenc}
\usepackage[utf8]{inputenc}
\usepackage{lmodern}
\usepackage[margin=2.2cm,top=2.0cm,bottom=2.3cm]{geometry}
\usepackage{microtype}
\usepackage{graphicx}
\usepackage{xcolor}
\usepackage{titlesec}
\usepackage{enumitem}
\usepackage{booktabs}
\usepackage{tabularx}
\usepackage{fancyhdr}
\usepackage{lastpage}
\usepackage{tcolorbox}
\usepackage{xurl}
\usepackage[hidelinks]{hyperref}

\tcbuselibrary{skins,breakable}

% Logos live in template/ and are shared by every issue -- never copy them into a
% dated folder. This path list resolves them from <YYYY-MM-DD>/ or from template/.
\graphicspath{{../../../template/}{./}{template/}}

%------------------------------ palette --------------------------------------
\definecolor{brandink}{HTML}{103868}    % KAUST navy -- headings, rules
\definecolor{brandaccent}{HTML}{0B7480} % KAUST teal (darkened to 5.6:1 on white)
\definecolor{brandsoft}{HTML}{EAF3F5}   % box fills
\definecolor{brandgrey}{HTML}{5A6672}   % meta text

\hypersetup{colorlinks=true,linkcolor=brandaccent,urlcolor=brandaccent,
            citecolor=brandaccent,pdfborder={0 0 0}}

%------------------------------ metadata -------------------------------------
% Issue number: one more than the highest \briefnumber in the existing dated folders.
% A REBUILD of a window that already has a folder keeps that folder's number -- the
% number identifies the window, not the run.
\newcommand{\briefnumber}{09}
% The window is seven days ending on \briefdate, INCLUSIVE -- so the span is 6 days
% wide (19--25, not 18--25). Check both endpoints against `date +%F` before shipping.
\newcommand{\briefweek}{14--20 October 2026}
\newcommand{\briefdate}{20 October 2026}            % date of publication
\newcommand{\briefauthor}{Ali Habibullah}          % prepared by
\newcommand{\briefaffil}{KAUST}                    % affiliation
\newcommand{\briefscope}{Papers, model releases, benchmarks and policy from the past seven days}

%------------------------------ section styling ------------------------------
\titleformat{\section}
  {\normalfont\sffamily\large\bfseries\color{brandink}}{\thesection.}{0.6em}{}
  [\vspace{-0.55em}{\color{brandaccent}\rule{\textwidth}{0.9pt}}]
\titlespacing*{\section}{0pt}{1.5ex plus .3ex}{0.9ex}

\titleformat{\subsection}[runin]
  {\normalfont\sffamily\bfseries\color{brandink}}{}{0pt}{}[\;---]
\titlespacing*{\subsection}{0pt}{1.25ex plus .2ex}{0.5em}

\setlist[itemize]{leftmargin=1.15em,itemsep=0.25em,topsep=0.35em,parsep=0pt}
\setlength{\parindent}{0pt}
\setlength{\parskip}{0.45em}

%------------------------------ header / footer ------------------------------
\pagestyle{fancy}
\fancyhf{}
\renewcommand{\headrulewidth}{0pt}
\renewcommand{\footrulewidth}{0.4pt}
\fancyfoot[L]{\footnotesize\color{brandgrey}Weekly AI Brief \#\briefnumber\ \textbullet\ \briefweek}
\fancyfoot[R]{\footnotesize\color{brandgrey}Page \thepage\ of \pageref{LastPage}}

%------------------------------ content macros -------------------------------
% Executive-summary box.  \tldr{ \item ... \item ... }
\newcommand{\tldr}[1]{%
  \begin{tcolorbox}[breakable,enhanced,colback=brandsoft,colframe=brandaccent,
      boxrule=0.8pt,left=10pt,right=10pt,top=8pt,bottom=8pt,arc=2pt,
      title={\sffamily\bfseries The week in five lines},fonttitle=\color{white}]
    \begin{itemize}[leftmargin=1.1em,itemsep=0.3em,topsep=0pt]#1\end{itemize}
  \end{tcolorbox}}

% One item of news.  \entry{Headline}{Who / venue}{URL}{source-number}
\newcommand{\entry}[4]{%
  \par\vspace{0.7em}%
  {\sffamily\bfseries\color{brandink}#1}\par
  {\footnotesize\color{brandgrey}#2 \textbullet\ \href{#3}{link} \textbullet\ \srcref{#4}}\par
  \vspace{0.15em}}

% "Why it matters" line -- one sentence, no hedging.
\newcommand{\why}[1]{\par{\small\textbf{Why it matters.} #1}\par}

% Numbered pointer into SOURCES.md / the source table.
\newcommand{\srcref}[1]{\textcolor{brandaccent}{[S#1]}}

% Numbers worth pulling out of a paragraph.
\newcommand{\stat}[2]{{\sffamily\bfseries\color{brandink}#1}~{\footnotesize\color{brandgrey}#2}}

%==============================================================================
\begin{document}

%------------------------------ title block ----------------------------------
\noindent
\begin{minipage}[c]{0.48\textwidth}
  \raggedright\includegraphics[height=1.4cm]{KAUST_Academy_Logo.png}
\end{minipage}\hfill
\begin{minipage}[c]{0.48\textwidth}
  \raggedleft\includegraphics[height=1.5cm]{KAUST_Logo.png}
\end{minipage}
\vspace{0.7em}

{\sffamily
 {\color{brandgrey}\footnotesize ISSUE \#\briefnumber\ \textbullet\ \briefdate}\par
 \vspace{0.15em}
 {\LARGE\bfseries\color{brandink}Weekly AI Brief}\par
 \vspace{0.25em}
 {\large\color{brandaccent}\briefweek}\par
 \vspace{0.35em}
 {\footnotesize\color{brandgrey}\briefscope\\
  Prepared by \briefauthor, \briefaffil}\par}
\vspace{0.3em}
{\color{brandink}\rule{\textwidth}{1.4pt}}
\vspace{0.6em}

%------------------------------ 1. executive summary -------------------------
\tldr{
  \item \textbf{Three large open-weight models shipped in three days.} Meta's Muse Spark Open (816B total, 42B active) scores 39 on Artificial Analysis's index, Reflection's Apache-2.0 Beam scores 30, and DeepSeek's MIT-licensed V4.1 Flash 1017 scores 45. \srcref{2}\,\srcref{5}\,\srcref{8}
  \item \textbf{A 132-run, 12,124 GPU-hour study finds no self-distillation method consistently beats plain GRPO} for agent RL, a week after two papers questioned on-policy distillation. \srcref{11}
  \item \textbf{Claude Sonnet 5.5 is now first on the Vals Index (70.85\%)} after a version change lifted every model. Counting refusal fallbacks as failures, Gemini 4 Argon would lead. \srcref{16}
  \item \textbf{The EU published its first harmonised standard under the AI Act,} EN 18286 for quality management systems, giving providers a presumption of conformity with Article~17. \srcref{22}
  \item \textbf{On 23 October OpenAI shuts down gpt-4, gpt-4-turbo, o1, o3-mini, o4-mini and gpt-image-1 in its API.} GPT-4o and GPT-4.1 follow on 14 January 2027. \srcref{25}
}

%------------------------------ 2. models & releases -------------------------
\section{Model and product releases}

\entry{Muse Spark Open: Meta releases the weights of a frontier-scale multimodal MoE}{Meta \textbullet\ open weights on Hugging Face (gated) \textbullet\ 15 October}{https://huggingface.co/blog/meta-muse-spark-open}{1}
Muse Spark Open is a mixture-of-experts model with about 816B total and 42B active parameters
(256 experts, 8 routed per token). It takes text, images, audio and video, outputs text, and has a 1M-token
context. Checkpoints come in BF16 and NVFP4. The repository is gated and tagged with a custom
(``other'') licence. Meta's own announcement and model card could not be read this week, so the
specifications here come from Hugging Face's launch post. Artificial Analysis scores it 39 (xhigh effort),
at 140 output tokens/s and \$1.25/\$4.25 per million input/output tokens through API providers. \srcref{2}
It is not Muse Wonder, last week's closed model, which scores 52.
\why{It is the strongest open-weight model from a US lab on Artificial Analysis's index, though the gate and custom licence mean teams should read the terms before building on it.}

\entry{Reflection Beam 1.0: the waitlisted preview becomes Apache-2.0 weights}{Reflection \textbullet\ open weights \textbullet\ 16 October}{https://huggingface.co/blog/reflection-beam}{3}
Beam-1.0-501B-A23B is a text-only MoE with 501B total and 23B active parameters, a 1M-token context, and
BF16, FP8 and NVFP4 checkpoints under Apache~2.0. Reflection's self-reported BF16 numbers include 81.6\% on
SWE-bench Verified, 80.1\% on Terminal-Bench~2.1 and 90.6\% on GPQA Diamond. \srcref{4} Independent tests are
much less flattering. Artificial Analysis scores Beam 30 on its index and 3.0\% on Terminal-Bench~4.0, its weakest
result. Its strongest is AutomationBench-AA, at 48.5\%. \srcref{5}
\why{It is the most permissively licensed model of the three, but the gap between Reflection's and Artificial Analysis's numbers means teams should run their own evaluation before adopting it.}

\entry{DeepSeek V4.1 Flash 1017: a point release that moves the open-weight lead}{DeepSeek \textbullet\ weights (MIT) and API \textbullet\ 17 October}{https://api-docs.deepseek.com/news/news261017}{6}
DeepSeek calls V4.1-Flash-1017 a ``minor update'' with shorter outputs. Its self-reported scores
rise from 36.8 to 42.6 on HLE without tools, from 31.2 to 42.3 on Terminal-Bench~4.0 and from 74.2 to 77.0 on DeepSWE~v1.1.
The model has a 1M-token context and MIT-licensed weights. \srcref{7} Artificial Analysis independently scores
it 45, second among open-weight models behind MiMo-V2.6-Pro (46) and level with GLM-5.3, at \$0.30/\$1.20 per million tokens. \srcref{8}
Vals places it eighth on its index at 61.26\%, with a cost of \$0.66 per test against \$3.61 for GPT-6.1 Sol. \srcref{16}
DeepSeek also opened its Search and Fetch APIs at \$0.75 per 1,000 calls. Calls from its harness,
Claude Code and Codex are free until billing starts on 31 October. \srcref{6}
\why{For cost-sensitive agent workloads it is one of the first open models to benchmark: near the top of the open-weight scores, MIT-licensed and cheap per token.}

\entry{Nemotron 3.5: NVIDIA's Super VL and a 24\,GB Bolt}{NVIDIA \textbullet\ open weights \textbullet\ 15 October}{https://blogs.nvidia.com/blog/nemotron-3-5-super-vl/}{9}
NVIDIA released Nemotron 3.5 Super VL and Nemotron 3.5 Bolt 74B-A9B, which is sized to run in 24\,GB. A 4\,GB Bolt E4B is due on
23 October. Artificial Analysis scores Super VL at 26, double Nemotron 3 Super, and Bolt at 21.4. It measured Super VL at
about 630 tokens/s with 100k-token prompts, on a pre-release endpoint. \srcref{10}
\why{Bolt is a credible local agent model for a single 24\,GB GPU.}

%------------------------------ 3. research ----------------------------------
\section{Research worth reading}

\entry{A reality check on self-distillation in agentic RL: none beats GRPO consistently}{Zhong et al.\ (incl.\ Chunhua Shen) \textbullet\ arXiv:2610.19789 \textbullet\ 14 October}{https://arxiv.org/abs/2610.19789}{11}
The authors re-run ten self-distillation methods against GRPO on ALFWorld, ScienceWorld, Search-QA and WebShop
at three model sizes. That is 132 training runs and 12,124 recorded H800 GPU-hours, with each evaluation repeated over eight rollout seeds.
No method beats GRPO across settings. Strong final scores can hide unstable training, and extra training
cost does not reliably buy better results. Code and a project page are public.
\why{Last week two papers questioned on-policy distillation; this one adds controlled evidence that the self-distillation variants also add little over the GRPO baseline.}

\entry{Skip the Compressor: give coding agents whole files, not compressed context}{Yang et al.\ (incl.\ ChengXiang Zhai), UIUC \textbullet\ arXiv:2610.21073 \textbullet\ 15 October}{https://arxiv.org/abs/2610.21073}{12}
At an equal per-issue token budget on 300 recent SWE-rebench issues, a simple selector (CQ2-CARE) that shows chosen
files verbatim resolves 43.7\% (GPT-5.6) and 46.7\% (Claude Opus~5). That is 6.3 points more than the best of
three prompt compressors, and 9.3 and 12.7 points more than the best retriever. It also edges out two trained methods, SweRank
and SWE-Pruner. The authors' explanation is that retrievers miss the files that need editing, and compressors alter the
text, so patches fail to apply. Code is not yet linked.
\why{If your agent compresses repository context to fit a budget, test spending that budget on choosing whole files instead.}

\entry{What RL brings to reasoning: about 4\% of tokens carry 90\% of the gain}{Zhang et al.\ \textbullet\ arXiv:2610.24663 \textbullet\ 17 October}{https://arxiv.org/abs/2610.24663}{13}
On 120 AIME problems with Qwen3-4B-Base, the authors take the RL-trained model's token only where the two models' greedy choices
differ, about 4\% of positions. That recovers 90.5\% of the RL model's pass@1 gain. Forcing the base model to open
with ``Alright'' alone recovers 35\%. At $k{=}256$ the sparse swap solves 67.5\% of problems, against 61.4\% for the RL model,
though at equal attempts rather than equal compute. The study covers one model size on AIME, plus three model pairs on MATH-500. Code is public.
\why{This adds evidence that reinforcement learning with verifiable rewards changes a few key decisions rather than adding new capability.}

\entry{COUNT: n-gram counts in host memory buy the effect of 57\% more compute}{Magnusson et al.\ (incl.\ Noah Smith, Pang Wei Koh) \textbullet\ arXiv:2610.20456 \textbullet\ 14 October}{https://arxiv.org/abs/2610.20456}{14}
COUNT stores precomputed corpus n-gram counts in ordinary host RAM and uses them to build an extra input embedding,
adding no learned parameters. In 1B-parameter pretraining at fixed compute, it improves validation loss as much as an
LM trained with 57\% more compute. Swapping in a count table from a domain-specific corpus at inference time adapts the model with no training.
Results are at 1B scale only. Code is linked from the paper.
\why{Host RAM sits mostly idle during training; this is a cheap way to use it and a training-free route to domain adaptation.}

%------------------------------ 4. benchmarks & evals ------------------------
\section{Benchmarks, evaluations and reproductions}

The Artificial Analysis (AA) index did not move at the top: it is still v4.3.2, with Opus~5.5 at 58. Vals moved to v2.2 on 19 October,
updating its finance, coding-migration and legal components. Every model in the table rose, so the gains reflect
the new version more than changes to the models. The order changed: Sonnet~5.5 passes Argon. Epoch re-estimated its scores by under 0.6 points and now
lists Argon and Haiku~5.5.

{\small
\begin{tabularx}{\textwidth}{@{}X r r r@{}}
\toprule
\textbf{Model (AA variant)} & \textbf{AA Index v4.3.2} & \textbf{Vals Index v2.2} & \textbf{Epoch ECI} \\
\midrule
Claude Opus 5.5 (max)          & 58 & 70.18\% (+3.21) & 166.82 \\
Claude Sonnet 5.5 (max)        & 56 & 70.85\% (+3.81) & 164.45 \\
Muse Wonder Full (max)         & 56 & 66.79\% (+2.22) & -- \\
Claude Fable 5.1 (max)         & 53 & -- (was 65.83\%) & 164.65 \\
GPT-6 Astra (max)              & 53 & 65.02\% (+1.89) & 166.41 \\
Gemini 4 Argon (high)          & 53 & 69.24\% (+0.34) & 164.56 \emph{new} \\
Muse Wonder (max)              & 52 & 64.88\% (+3.82) & -- \\
GPT-6.1 Sol (max)              & 52 & 63.94\% (+2.79) & 166.32 \\
DeepSeek V4.1 Flash 1017 (max) \emph{new} & 45 & 61.26\% & -- \\
Claude Haiku 5.5 (max)         & 43 & 58.71\% & 158.98 \emph{new} \\
\bottomrule
\end{tabularx}}
\par\vspace{0.2em}
{\footnotesize\color{brandgrey}Sources: AA leaderboard, accessed 20 Oct \srcref{15}; Vals Index v2.2, updated 19 Oct \srcref{16};
Epoch ECI data, accessed 20 Oct \srcref{17}. Brackets: change against v2.1 in issue \#08. ``--'' = not listed. Anthropic rows on AA
and Vals include fallback-model answers after refusals.}
\vspace{0.3em}

\textbf{Is the movement real?} Sonnet~5.5's lead on Vals is not solid. Vals itself notes that counting
refusal-fallback tasks as failures drops Sonnet to 69.03\% and Opus to 64.16\%, which would put Argon first. Fable~5.1
is no longer in the Vals table, and the page gives no reason. Epoch's top five are within 2.4 points, inside each other's confidence
intervals (Opus 163.54--170.99). Two new evaluations are worth noting:

\begin{itemize}
  \item \textbf{AA-Robotics v0.1} gives models a 6-axis robot arm through a coding-agent harness, across eight tabletop tasks. GPT-6 Astra scores 69.4, GPT-6.1 Sol
  64.6, Argon 55.3, Opus~5.5 48.8 and Fable~5.1 46.3. Astra costs \$16.36 per task, more than eight times Sol's \$1.93, for about
  four points more. These are early results. \srcref{18}
  \item \textbf{HealthBench Professional on AA} (525 clinician conversations): after adjusting scores for answer length, GPT-6 Astra leads at 66.0\%.
  Opus~5.5 has the top raw score, 71.4\%, but drops to 59.9\% after the adjustment. On the broader HealthBench, Sonnet~5.5 leads at 62.1\%. \srcref{19}
\end{itemize}

%------------------------------ 5. industry & policy -------------------------
\section{Industry, funding and policy}

\begin{itemize}
  \item \textbf{Tenstorrent raises \$1.5B at a \$6.2B valuation (19 October).} The Series E was led by Trace Capital and
  Innovation X, with ARK Invest, Fidelity, Morgan Stanley, T.~Rowe Price, Sandisk and XTX Markets among the participants. The
  money goes to scaling up Galaxy servers and to the RISC-V CPU and AI core designs it licenses to chipmakers. \srcref{20}
  \item \textbf{Anthropic commits 5\% of its shares to a new Public Benefit Fund (18 October).} The commitment is subject to stockholder approval, and the
  board may add 2\% for each 50\% rise in company value above a baseline it has not yet disclosed. Including the co-founders' 80\% equity pledge and
  the employee donation match, Anthropic says about 20\% of its shares will go to the public good. The fund is an LLC, not a charity, and none of its equity may be used for lobbying or
  political giving. \srcref{21}
  \item \textbf{The EU publishes its first harmonised standard under the AI Act (19 October).} EN 18286:2026, written by CEN and CENELEC and referenced
  via Implementing Decision (EU) 2026/2274, covers quality management systems. Providers of high-risk AI systems that implement it are presumed to
  conform with the Article~17 requirements it covers. The Commission says more harmonised standards will follow. \srcref{22}
  \item \textbf{US agencies publish incident practices for super intelligence training (16 October).} NIST's Center for Advancing Innovation and Standards for Super Intelligence (CAISSI),
  the NSA's Super Intelligence Security Center and CISA issued ten voluntary practices across five stages of incident response. They target incidents in which AI agents take
  significant actions nobody instructed and that break operator controls. \srcref{23}
  \item \textbf{Iowa sues OpenAI under its Consumer Fraud Act (14 October).} The Iowa attorney general's release says Arkansas, Arizona, Idaho and
  Montana have filed their own suits in state courts. The attorneys general of Kansas, Nebraska and West Virginia joined the announcement. The
  release gives no damages figure. \srcref{24}
\end{itemize}

%------------------------------ 6. what to watch -----------------------------
\section{What to watch next week}

\begin{itemize}
  \item \textbf{23 October: OpenAI API shutdowns.} gpt-4, gpt-4-turbo, gpt-3.5-turbo, o1, o1-pro, o3-mini, o4-mini, gpt-image-1
  and several fine-tune bases are removed. OpenAI's suggested replacements are mainly gpt-5.6-sol and gpt-5.6-terra. \srcref{25}
  \item \textbf{23 October: Nemotron 3.5 Bolt E4B,} NVIDIA's 4\,GB edge model, is due. Watch for independent scores. \srcref{9}
  \item \textbf{December: the EU AI Act's ninth prohibition takes effect.} It covers systems that generate non-consensual intimate content or
  child sexual abuse material, such as ``nudification'' apps, and was added through the AI Omnibus. The Commission's page gives the month, not the day. \srcref{26}
\end{itemize}

%------------------------------ 7. sources -----------------------------------
\clearpage
\section{Sources}
{\footnotesize
\begin{tabularx}{\textwidth}{@{}l X@{}}
\toprule
\textbf{\#} & \textbf{Reference} \\
\midrule
S1  & Hugging Face. \emph{Welcome Muse Spark Open}. 15 Oct 2026. \url{https://huggingface.co/blog/meta-muse-spark-open} \\
S2  & Artificial Analysis. \emph{Muse Spark Open}. Accessed 20 Oct 2026. \url{https://artificialanalysis.ai/models/muse-spark-open} \\
S3  & Hugging Face. \emph{Welcome Beam: a 501B open model}. 16 Oct 2026. \url{https://huggingface.co/blog/reflection-beam} \\
S4  & Reflection. \emph{Beam-1.0-501B-A23B} model card. Hugging Face. \url{https://huggingface.co/reflection/Beam-1.0-501B-A23B} \\
S5  & Artificial Analysis. \emph{Reflection Beam}. 16 Oct 2026. \url{https://artificialanalysis.ai/articles/reflection-beam} \\
S6  & DeepSeek. \emph{DeepSeek-V4.1-Flash-1017 Release}. 17 Oct 2026. \url{https://api-docs.deepseek.com/news/news261017} \\
S7  & DeepSeek. \emph{DeepSeek-V4.1-Flash-1017} model card. Hugging Face. \url{https://huggingface.co/deepseek-ai/DeepSeek-V4.1-Flash-1017} \\
S8  & Artificial Analysis. \emph{DeepSeek V4.1 Flash 1017}. Accessed 20 Oct 2026. \url{https://artificialanalysis.ai/models/deepseek-v4-1-flash-1017} \\
S9  & NVIDIA. \emph{Nemotron 3.5 Super VL}. 15 Oct 2026. \url{https://blogs.nvidia.com/blog/nemotron-3-5-super-vl/} \\
S10 & Artificial Analysis. \emph{NVIDIA Releases Nemotron 3.5 Super and Nemotron 3.5 Bolt}. 15 Oct 2026. \url{https://artificialanalysis.ai/articles/nemotron-3-5-super-and-bolt} \\
S11 & Zhong et al. \emph{More Methods, Better Agents? A Reality Check on Self-Distillation in Agentic RL}. arXiv:2610.19789, 14 Oct 2026. \url{https://arxiv.org/abs/2610.19789} \\
S12 & Yang et al. \emph{Skip the Compressor}. arXiv:2610.21073, 15 Oct 2026. \url{https://arxiv.org/abs/2610.21073} \\
S13 & Zhang et al. \emph{What RL Brings to Reasoning}. arXiv:2610.24663, 17 Oct 2026. \url{https://arxiv.org/abs/2610.24663} \\
S14 & Magnusson et al. \emph{The COUNT}. arXiv:2610.20456, 14 Oct 2026. \url{https://arxiv.org/abs/2610.20456} \\
S15 & Artificial Analysis. \emph{Models leaderboard}. Accessed 20 Oct 2026. \url{https://artificialanalysis.ai/leaderboards/models} \\
S16 & Vals AI. \emph{Vals Index} v2.2. Updated 19 Oct 2026. \url{https://www.vals.ai/benchmarks/vals_index} \\
S17 & Epoch AI. \emph{ECI scores} (CSV). Accessed 20 Oct 2026. \url{https://epoch.ai/data/eci_scores.csv} \\
S18 & Artificial Analysis. \emph{AA-Robotics v0.1}. 14 Oct 2026. \url{https://artificialanalysis.ai/articles/aa-robotics} \\
S19 & Artificial Analysis. \emph{HealthBench Professional}. 19 Oct 2026. \url{https://artificialanalysis.ai/articles/healthbench-professional-health-conversations} \\
S20 & Tenstorrent. \emph{Tenstorrent Announces a \$1.5B Series E at a \$6.2B valuation}. 19 Oct 2026. \url{https://tenstorrent.com/newsroom/tenstorrent-announces-series-e} \\
S21 & Anthropic. \emph{Investing in the public benefit}. 18 Oct 2026. \url{https://www.anthropic.com/news/public-benefit-fund} \\
S22 & European Commission. \emph{First harmonised standard supporting the AI Act published in the Official Journal}. 19 Oct 2026. \url{https://digital-strategy.ec.europa.eu/en/library/first-harmonised-standard-supporting-ai-act-published-official-journal} \\
S23 & NIST. \emph{Key Practices for Incident Readiness and Response in Super Intelligence Training and Evaluation}. 16 Oct 2026. \url{https://www.nist.gov/news-events/news/2026/10/key-practices-incident-readiness-and-response-super-intelligence-training} \\
S24 & Iowa Attorney General. \emph{Lawsuit against OpenAI}. 14 Oct 2026. \url{https://www.iowaattorneygeneral.gov/newsroom/attorney-general-brenna-bird-announces-lawsuit-against-openai-for-deceiving-iowa-parents-and-iowans} \\
S25 & OpenAI. \emph{Deprecations}. Accessed 20 Oct 2026. \url{https://developers.openai.com/api/docs/deprecations} \\
S26 & European Commission. \emph{AI Act: regulatory framework}. Accessed 20 Oct 2026. \url{https://digital-strategy.ec.europa.eu/en/policies/regulatory-framework-ai} \\
\bottomrule
\end{tabularx}}

\vspace{0.8em}
{\footnotesize\color{brandgrey}
Compiled \briefdate. Every figure in this brief was read from the linked primary
source on that date. Corrections: \texttt{\briefauthor}.}

\end{document}
